% This document was generated by an LLM.

\documentclass{essay}

\title{Higher Implicit Derivatives}
\author{}
\date{}

\begin{document}

\maketitle

\section{Setup}

Let $F(x,y)$ be a smooth (say $C^\infty$) function of two variables, and consider the level
set
\[ F(x,y) = C \]
for some constant $C$. Near any point where $F_y \neq 0$, the Implicit Function Theorem
guarantees that this equation defines $y$ as a function of $x$, say $y = y(x)$, at least
locally. We are interested in computing the derivatives $y'$, $y''$, $y'''$, \dots{} directly
from $F$ and its partial derivatives, without ever solving for $y(x)$ explicitly.

The method is always the same: differentiate the identity $F(x,y(x)) = C$ with respect to
$x$, using the chain rule, and then solve algebraically for whichever derivative of $y$ is
new. Each differentiation introduces one new unknown ($y'$, then $y''$, then $y'''$, \dots{})
and one new equation, so the system is always solvable as long as $F_y \neq 0$.

Throughout, subscripts on $F$ denote partial derivatives: $F_x = \partial F/\partial x$,
$F_{xy} = \partial^2 F/\partial x\, \partial y$, and so on. We assume enough smoothness that
mixed partials commute (Clairaut's theorem), e.g.\ $F_{xy} = F_{yx}$.

\section{First derivative}

Differentiating $F(x,y(x)) = C$ once with respect to $x$:
\[
  F_x + F_y\, y' = 0
  \quad\Longrightarrow\quad
  y' = -\frac{F_x}{F_y}.
\]
This is the standard implicit differentiation formula. Nothing new yet, but it sets the
pattern: every subsequent step is ``differentiate the previous equation again, then solve
for the new unknown.''

\section{Second derivative}

Differentiate $F_x + F_y y' = 0$ once more with respect to $x$. The key point is that $F_x$
and $F_y$ are themselves functions of $x$ and $y$, so differentiating them requires the
chain rule again:
\[
  \frac{d}{dx}F_x = F_{xx} + F_{xy}\,y', \qquad
  \frac{d}{dx}F_y = F_{yx} + F_{yy}\,y'.
\]
Combined with the product rule on the term $F_y y'$, this gives
\[
  F_{xx} + F_{xy}y' + \big(F_{yx} + F_{yy}y'\big)y' + F_y y'' = 0.
\]
Using $F_{xy} = F_{yx}$ and collecting terms:
\[
  F_{xx} + 2F_{xy}y' + F_{yy}(y')^2 + F_y y'' = 0.
\]
Solving for $y''$:
\[
  y'' = -\frac{F_{xx} + 2F_{xy}y' + F_{yy}(y')^2}{F_y}.
\]
This expression still contains $y'$; substituting $y' = -F_x/F_y$ and clearing denominators
gives the closed form
\[
  \boxed{\,y'' = -\dfrac{F_{xx}F_y^2 - 2F_{xy}F_xF_y + F_{yy}F_x^2}{F_y^3}\,}.
\]

\paragraph{Connection to the bordered Hessian.} The numerator is (up to sign) the
determinant
\[
  \begin{vmatrix}
    F_{xx} & F_{xy} & F_x \\
    F_{xy} & F_{yy} & F_y \\
    F_x    & F_y    & 0
  \end{vmatrix}
  = F_y^2 F_{xx} - 2F_xF_yF_{xy} + F_x^2F_{yy},
\]
which is exactly the numerator above. This matrix (the Hessian of $F$ bordered by its
gradient) is the same object that appears in second-derivative tests for constrained
optimization and in formulas for the curvature of a level curve. The recurring appearance
of this determinant is not a coincidence: $y''$, curvature, and the second-order behavior of
Lagrange-multiplier problems are all measuring ``how much $F$ curves relative to how steeply
it is changing,'' which is precisely what a bordered Hessian encodes.

\section{Third derivative}

The same procedure continues. Differentiate
\[
  F_{xx} + 2F_{xy}y' + F_{yy}(y')^2 + F_y y'' = 0
\]
once more with respect to $x$, remembering that every partial of $F$ again picks up a
$(\,\cdot\,)_y \cdot y'$ term via the chain rule, and that $2F_{xy}y'$, $F_{yy}(y')^2$, and
$F_y y''$ each need the product rule. Carrying this out and simplifying (using $F_{xy} =
F_{yx}$ and $F_{xyy} = F_{yxy} = F_{yyx}$, etc.) yields
\[
  F_{xxx} + 3F_{xxy}y' + 3F_{xyy}(y')^2 + F_{yyy}(y')^3 + 3F_{xy}y'' + 3F_{yy}y'y'' + F_y y''' = 0,
\]
so that
\[
  y''' = -\frac{F_{xxx} + 3F_{xxy}y' + 3F_{xyy}(y')^2 + F_{yyy}(y')^3 + 3F_{xy}y'' +
  3F_{yy}y'y''}{F_y}.
\]
As with $y''$, one can substitute the already-known expressions for $y'$ and $y''$ to get a
formula purely in terms of $F$'s partial derivatives, but the resulting expression is long
and rarely worth memorizing. The binomial-looking coefficients ($1,3,3,1$) are not an
accident: they come from repeatedly applying the chain and product rules to $F$ composed
with $(x, y(x))$, which is structurally close to a Faà di Bruno / Leibniz expansion.

\section{The general pattern}

Every order follows the same three-step recipe:
\begin{enumerate}[label=(\arabic*)]
  \item Start from the equation obtained at the previous order (at order 1, this is
    $F(x,y) = C$ itself).
  \item Differentiate the whole equation with respect to $x$, applying the chain rule to every
    term involving $F$ or its partials (since they depend on $y$, which depends on $x$), and the
    product rule to every term already containing a derivative of $y$.
  \item Collect terms and solve algebraically for the new, highest-order derivative of $y$
    that appears. It will always show up linearly, multiplied by $F_y$, because each
    differentiation of $F(x, y(x))$ with respect to $x$ produces exactly one term of the form
    $F_y \cdot (\text{new derivative})$ from the chain rule applied to the $y$-dependence.
\end{enumerate}
This is why the requirement $F_y \neq 0$ is exactly what is needed to solve for $y'$, $y''$,
$y'''$, and so on: at every order, $F_y$ is the coefficient of the new unknown.

\section{Worked examples}

\subsection*{Example 1: the unit circle}

Let $F(x,y) = x^2 + y^2$ and $C = 1$, i.e.\ the unit circle $x^2 + y^2 = 1$.

Partials: $F_x = 2x$, $F_y = 2y$, $F_{xx} = 2$, $F_{yy} = 2$, $F_{xy} = 0$.

First derivative:
\[
  y' = -\frac{2x}{2y} = -\frac{x}{y}.
\]

Second derivative, using the boxed formula:
\[
  y'' = -\frac{F_{xx}F_y^2 - 2F_{xy}F_xF_y + F_{yy}F_x^2}{F_y^3}
  = -\frac{2(2y)^2 - 0 + 2(2x)^2}{(2y)^3}
  = -\frac{8y^2 + 8x^2}{8y^3}
  = -\frac{x^2+y^2}{y^3}.
\]
Since $x^2 + y^2 = 1$ on the circle, this simplifies to
\[
  y'' = -\frac{1}{y^3}.
\]
This matches the direct computation from $y = \sqrt{1-x^2}$, and it has a clean geometric
meaning: for the unit circle, curvature is constant, and $y''$ reflects that constancy once
restricted to the curve.

\subsection*{Example 2: a curve with a mixed term}

Let $F(x,y) = x^2 + xy + y^2$ and $C = 3$ (an ellipse, since the quadratic form is positive
definite).

Partials: $F_x = 2x + y$, $F_y = x + 2y$, $F_{xx} = 2$, $F_{yy} = 2$, $F_{xy} = 1$.

First derivative:
\[
  y' = -\frac{2x+y}{x+2y}.
\]

Second derivative:
\[
  y'' = -\frac{2(x+2y)^2 - 2(1)(2x+y)(x+2y) + 2(2x+y)^2}{(x+2y)^3}.
\]
Expanding the numerator:
\[
  2(x+2y)^2 = 2x^2 + 8xy + 8y^2, \qquad
  2(2x+y)(x+2y) = 4x^2 + 10xy + 4y^2,
\]
\[
  2(2x+y)^2 = 8x^2 + 8xy + 2y^2,
\]
so the numerator is
\[
  (2x^2+8xy+8y^2) - (4x^2+10xy+4y^2) + (8x^2+8xy+2y^2) = 6x^2 + 6xy + 6y^2 = 6(x^2+xy+y^2).
\]
Using $x^2 + xy + y^2 = 3$ on the curve, the numerator is $18$, so
\[
  y'' = -\frac{18}{(x+2y)^3}.
\]
This example shows the mixed-partial term $F_{xy}$ doing real work: without the $-2F_{xy}F_xF_y$
term, the cancellation down to a multiple of $F$ itself would not happen.

\subsection*{Example 3: a transcendental curve}

Let $F(x,y) = \sin(x) + e^{y} - xy$ and $C = 1$, i.e.\ $\sin x + e^y - xy = 1$.

Partials: $F_x = \cos x - y$, $F_y = e^y - x$, $F_{xx} = -\sin x$, $F_{yy} = e^y$,
$F_{xy} = -1$.

First derivative:
\[
  y' = -\frac{\cos x - y}{e^y - x}.
\]

Second derivative, directly from the boxed formula (no further simplification is available
here, since the curve is not defined by a homogeneous polynomial):
\[
  y'' = -\frac{(-\sin x)(e^y-x)^2 - 2(-1)(\cos x - y)(e^y - x) + (e^y)(\cos x - y)^2}
  {(e^y - x)^3}.
\]
This can be evaluated at any specific point $(x_0, y_0)$ on the curve by plugging in numbers;
it illustrates that the method works identically whether or not $F$ is polynomial, since
nothing in the derivation used polynomiality — only differentiability.

\section{Self-check questions}

\begin{enumerate}[label=\arabic*.]
  \item Why does the coefficient of the newest derivative of $y$ (i.e.\ $y'$ in the first
    equation, $y''$ in the second, and so on) always turn out to be exactly $F_y$, at every
    order?
  \item Derive $y''$ from scratch for $F(x,y) = x^3 - 3xy + y^3$, $C = 0$ (the folium of
    Descartes), without looking at the boxed formula. Then check your answer against the boxed
    formula.
  \item The boxed formula for $y''$ can be written as a $3\times 3$ determinant divided by
    $F_y^3$. What goes wrong (in terms of that determinant, or in terms of the original
    derivation) if $F_y = 0$ at a point on the curve? What does this correspond to geometrically?
  \item In the derivation of $y'''$, coefficients $1, 3, 3, 1$ appeared. Where exactly do they
    come from? (Hint: track how many times the product rule is invoked on terms that already
      contain a $y'$ or $y''$ factor, versus how many times the chain rule is invoked on partials
    of $F$.)
  \item Suppose $F(x,y) = f(x) + g(y)$ for some single-variable functions $f, g$ (i.e.\ $F$ is
    ``separable''). Show that the general boxed formula for $y''$ collapses to something much
    simpler. Can you explain why separability should make the computation easier, structurally?
  \item If $F(x,y) = C$ implicitly defines $x$ as a function of $y$ instead (i.e.\ you solve
    for $x(y)$), what is the analogous formula for $x''(y)$? What symmetry do you expect between
    the $y''(x)$ and $x''(y)$ formulas, and does the boxed formula actually have that symmetry?
  \item For Example 1 (the unit circle), the formula for $y''$ simplified using the constraint
    $x^2+y^2=1$. Explain in general terms why it is valid to substitute the original constraint
    $F(x,y)=C$ back into a formula for $y''$, even though $y''$ was derived by differentiating,
    not by using the constraint directly.
  \item Explain, in your own words and without more computation, why $y'''$ needs $F_y \neq 0$
    at exactly the same points where $y''$ and $y'$ do — i.e.\ why no new nondegeneracy condition
    appears at higher orders.
\end{enumerate}

\end{document}
